\subsection{Sums of powers}

We again write \(\mathbf{X} = (\mathbf{X}_{i})_{i \in I}\) for a family of indeterminates. The symmetric formal power series \(s_{k}\) are defined as before by

\[
s _ {k} = \sum_ {\mathrm{H} \in \mathfrak {P} _ {k}} \prod_ {i \in \mathrm{H}} X _ {i} \quad (k \geqslant 1),\tag{16}
\]

where \(P_{k}\) is the set of all k-element subsets of I. We also write

\[
p _ {k} = \sum_ {i \in I} X _ {i} ^ {k} \quad (k \geqslant 1).\tag{17}
\]

This is a symmetric formal power series which is homogeneous of degree k.

Lemma 4 (« Newton's relations »). --- For every integer \(d \geqslant 1\) we have

\[
p _ {d} = \sum_ {k = 1} ^ {d - 1} (- 1) ^ {k - 1} s _ {k} p _ {d - k} + (- 1) ^ {d + 1} d s _ {d}.\tag{18}
\]

Let us define a continuous derivation \(\Delta\) in \(\mathbf{A}[[\mathbf{X}]]\) by \(\Delta(u) = \sum_{n \geq 0} nu_n\) , where for every \(u\) in \(\mathbf{A}[[\mathbf{X}]]\) , \(u_n\) is the homogeneous component of degree \(n\) of \(u\) . By (16) and Prop. 2 of IV, p.~27 we have

\[
1 + \sum_ {k \geqslant 1} s _ {k} = \prod_ {i \in I} (1 + X _ {i}).\tag{19}
\]

By Prop. 9 of IV, p.~34 we thus have

\[
\left(\sum_ {k \geqslant 1} k s _ {k}\right) \cdot \left(1 + \sum_ {k \geqslant 1} s _ {k}\right) ^ {- 1} = \sum_ {i \in I} \Delta \left(\mathrm{X} _ {i}\right) / \left(1 + \mathrm{X} _ {i}\right).\tag{20}
\]

We have \(\Delta(X_i) = X_i\) and \(X_i / (1 + X_i) = \sum_{k \geq 1} (-1)^{k-1} X_i^k\) . The right-hand side of (20) is thus equal to \(\sum_{k \geq 1} (-1)^{k-1} p_k\) . Hence by (20) we have

\[
\sum_ {k \geqslant 1} k s _ {k} = \left(1 + \sum_ {k \geqslant 1} s _ {k}\right) \cdot \left(\sum_ {k \geqslant 1} (- 1) ^ {k - 1} p _ {k}\right)
\]

and now Lemma 4 follows by a comparison of homogeneous components of degree \(d\) .

Remark. --- With the notation of the above proof we have

\[
\Delta u = \sum_ {i \in I} X _ {i} \cdot D _ {i} (u)
\]

(IV, p.~33, Cor. 1). In other words, Euler's relation (IV, p.~8, Prop 6) extends to formal power series: if \(u \in \mathbf{A}[\mathbf{X}]\) is homogeneous of degree \(n\) , then

\[
n \cdot u = \sum_ {i \in I} X _ {i} \cdot D _ {i} (u).\tag{21}
\]

When I is finite, consisting of \(n\) elements, we have \(s_k = 0\) for \(k > n\) . Then Newton's relations may be written

\[
\begin{array}{l} p _ {2} = s _ {1} p _ {1} - 2 s _ {2} \\ p _ {3} = s _ {1} p _ {2} - s _ {2} p _ {1} + 3 s _ {3} \\ \dots \dots \\ p _ {n - 1} = s _ {1} p _ {n - 2} - s _ {2} p _ {n - 3} + \dots + (- 1) ^ {n - 1} s _ {n - 2} p _ {1} + (- 1) ^ {n} (n - 1) s _ {n - 1} \\ p _ {n} = s _ {1} p _ {n - 1} - s _ {2} p _ {n - 2} + \dots + (- 1) ^ {n} s _ {n - 1} p _ {1} + (- 1) ^ {n + 1} n s _ {n} \end{array}
\]

and

\[
p _ {k} = s _ {1} p _ {k - 1} - s _ {2} p _ {k - 2} + \dots + (- 1) ^ {n + 1} s _ {n} p _ {k - n} \quad (\text { for } k > n).\tag{22}
\]

The first n of the above relations hold for any I; we find for example

\[
p _ {1} = s _ {1}, \quad p _ {2} = s _ {1} ^ {2} - 2 s _ {2}, \quad p _ {3} = s _ {1} ^ {3} - 3 s _ {1} s _ {2} + 3 s _ {3}.
\]

More generally, let \(\mathbf{S} = (\mathbf{S}_n)_{n\geqslant 1}\) be a family of indeterminates. Let us define the polynomials \(\mathbf{P}_d\in \mathbf{Z}[\mathbf{S}_1,\dots ,\mathbf{S}_d]\) recursively by \(\mathbf{P}_1 = \mathbf{S}_1\) and

\[
\mathrm{P} _ {d} = \sum_ {k = 1} ^ {d - 1} (- 1) ^ {k - 1} \mathrm{S} _ {k} \mathrm{P} _ {d - k} + (- 1) ^ {d + 1} d \mathrm{S} _ {d} \quad (d \geqslant 2).
\]

Then we have the « universal formulae » \(p_d = \mathrm{P}_d(s_1, \ldots, s_d)\) holding over any ring A and family X of indeterminates.

Let \(\mathbf{P} = (\mathrm{P}_k)_{k\geqslant 1}\) be a sequence of indeterminates. Since \(p_k\) is homogeneous of degree \(k\) in \(\mathbf{A}[[\mathbf{X}]]\) , there exists precisely one continuous A-algebra homomorphism

\[
\lambda_ {I}: \mathbf {A} [ [ \mathbf {P} ] ] \rightarrow \mathbf {A} [ [ \mathbf {X} ] ] ^ {\text { sym }}
\]

such that \(\lambda_{I}(\mathbf{P}_{k}) = p_{k}\) for every integer \(k \geqslant 1\) (IV, p.~28). If we assign to \(\mathbf{P}_{k}\) the weight \(k\) , then \(\lambda_{I}\) transforms an isobaric polynomial of weight \(n\) in the \(\mathbf{P}_{k}\) into a formal power series which is homogeneous of degree \(n\) in the \(\mathbf{X}_{i}\) .

PROPOSITION 3. --- a) If I is a finite set with n elements and \(n! \cdot 1\) is invertible in A, then \(\lambda_{I}\) induces a bicontinuous isomorphism of \(A[[P_{1}, ..., P_{n}]]\) onto \(A[[X]]^{sym}\) .

\begin{enumerate}
\def\labelenumi{\alph{enumi})}
\setcounter{enumi}{1}
\tightlist
\item
  If I is infinite and A is a Q-algebra, then \(\lambda_{I}\) is a bicontinuous isomorphism of A{[}{[}P{]}{]} onto A{[}{[}X{]}{]} \(^{\text{sym}}\) .
\end{enumerate}

We shall treat the case a) only, the case b) being quite similar.

By Th. 2 (IV, p.~68) we can identify A{[}{[}X{]}{]} \(^{sym}\) with the algebra of formal power series A{[}{[}S \(_1\) , \ldots, S \(_n\) {]}{]}, with S \(_k\) corresponding to s \(_k\) . By Lemma 4 of IV, p.~70, there exist formal power series g \(_1\) , \ldots, g \(_n\) of order \(\geqslant 2\) in the indeterminates s \(_1\) , \ldots, s \(_n\) such that

\[
p _ {k} = (- 1) ^ {k + 1} k s _ {k} + g _ {k} (s _ {1}, \dots , s _ {n}) \quad (1 \leqslant k \leqslant n).
\]

Since \(k!\) . 1 is invertible in A, Lemma 2 of IV, p.~35 proves the existence of an automorphism T of the topological A-algebra \(\mathbf{A}[[\mathbf{X}]]^{\mathrm{sym}}\) which maps \(s_k\) to \(p_k\) for \(1 \leqslant k \leqslant n\) . Now Prop. 3, a) is an immediate consequence.

COROLLARY. --- Let \(\xi_1, \ldots, \xi_n, \eta_1, \ldots, \eta_n\) be elements of A and suppose that A is an integral domain.

\begin{enumerate}
\def\labelenumi{\alph{enumi})}
\item
  If \(s_k(\xi_1, \ldots, \xi_n) = s_k(\eta_1, \ldots, \eta_n)\) for \(1 \leqslant k \leqslant n\) , then there exists a permutation \(\sigma \in \mathfrak{S}_n\) such that \(\eta_i = \xi_{\sigma(i)}\) for \(1 \leqslant i \leqslant n\) .
\item
  Suppose that \(n!\) . \(1 \neq 0\) in A and
\end{enumerate}

\[
\xi_ {1} ^ {k} + \dots + \xi_ {n} ^ {k} = \eta_ {1} ^ {k} + \dots + \eta_ {n} ^ {k}\tag{23}
\]

for \(1 \leqslant k \leqslant n\) . Then there exists a permutation \(\sigma \in \mathfrak{S}_n\) such that \(\eta_i = \xi_{\sigma(i)}\) for \(1 \leqslant i \leqslant n\) .

Under the hypotheses \(a)\) we have \(\prod_{i=1}^{n} (\mathbf{X} - \xi_i) = \prod_{i=1}^{n} (\mathbf{X} - \eta_i)\) . If we substitute \(\eta_n\) for \(\mathbf{X}\) , we find that \(\prod_{i=1}^{n} (\eta_n - \xi_i) = 0\) and since \(\mathbf{A}\) is an integral domain, there is an integer \(\sigma(n)\) such that \(1 \leqslant \sigma(n) \leqslant n\) and \(\eta_n = \xi_{\sigma(n)}\) . Now the assertion \(a)\) follows easily by induction because \(\mathbf{A}[\mathbf{X}]\) is an integral domain.

Under the hypotheses of \(b)\) there exist by Prop. 3 polynomials \(\Pi_1, \ldots, \Pi_n\) in \(n\) indeterminates, with coefficients in the field of fractions of \(A\) such that \(s_k = \Pi_k(p_1, \ldots, p_n)\) for \(1 \leqslant k \leqslant n\) . Now the relation (23) implies

\[
s _ {k} (\xi_ {1}, \dots , \xi_ {n}) = s _ {k} (\eta_ {1}, \dots , \eta_ {n})
\]

for \(1 \leqslant k \leqslant n\) and so \(b)\) follows from \(a)\) .

